\chapter{Literature Review}\label{chapter:lit_review}

Mediation analysis methods have historically been implemented in the social sciences 
within the linear structural equation modeling framework~\citep{wright1934,baronkenny1986}.
This framework even allows one to compute path-specific mediation effects given an 
arbitrary number of causally related mediators. However, this framework also 
imposes implausibly stringent assumptions on the functional form 
and potential interactions between treatment, outcome, and mediators~\citep{pearl2001}. 

More recently, the development of the potential outcomes framework for causal 
inference~\citep{rubin1974} has led to more flexible formulations of direct and indirect 
effects for mediation analysis. In particular,~\cite{robinsgreenland1992} leveraged 
the potential outcomes framework to formulate natural direct and indirect effects as 
counterfactual quantities that are not bound by the unrealistic relationships imposed by 
classical mediation analysis. This work was followed by another foundational paper that 
defined natural effects within the nonparametric structural equation modeling 
framework~\cite{pearl2001}. 

Despite the conceptual elegance of natural effects, concerns have been raised about the 
stringent assumptions needed to identify natural 
effects~\citep{vanderweelevansteelandt2009,imaietal2010b,vanderweele2015}. 
Of particular concern is the ”cross-world independence” assumption, 
which is not empirically testable; in fact, there does not even
exist a hypothetical experiment that could falsify such an assumption~\citep{robinsrichardson2010}. 
Additionally, the natural effects framework does not generalize well to situations in which
there is either a treatment-induced mediator-outcome confounder or where there is more than one
mediator of interest. Situations like these are common in the 
social and biomedical sciences~\citep{vansteelandtdaniel2017}, 
but natural effects cannot be identified here even if the cross-world identification assumption 
otherwise holds~\citep{avinetal2005,vanderweeleetal2014}. Finally, there is debate among 
researchers about the policy relevance of natural indirect effects because of the 
difficulty of conceptualizing a hypothetical intervention that could map to the natural 
indirect effect~\citep{vanderweelevansteelandt2009}.

In response to these issues,~\cite{vanderweele2014} proposed randomized interventional direct
and indirect effects as alternative, policy relevant effects that do not require untestable cross-world
independence assumptions for identification. These effects represent the effect of a joint stochastic
intervention on the treatment and the mediator that could hypothetically be implemented in an
experiment (see Section~\ref{chapter:interventional_effects}). These effects are similar to a 
more general version of direct and indirect effects defined without the use of counterfactuals 
in~\cite{didelez2006},~\cite{geneletti2007}, and~\cite{vanderlaanpetersen2008}. 

There have been two main strands of the theoretical literature on mediation analysis in recent years. 
The first strand remains committed to the value of natural effects estimation and focuses on developing
sensitivity analyses and partial identification bounds for natural effects that allow one to relax the 
stringent cross-world independence assumption. Work on sensitivity analyses for testing unmeasured 
confounding include work by~\cite{imaietal2010a} and~\cite{smithvanderweele2019}. Work on developing
bounds for natural effects has been done 
by~\cite{sjolander2009},~\cite{robinsrichardson2010},~\cite{tchetgentchetgenphiri2014}, 
and~\cite{milesetal2017}. There has also been interesting work done on developing
experiments that could, with a strong set of auxilliary assumptions, identify natural 
effects~\cite{robinsrichardson2010, imaietal2013}. While these experiments still impose strong 
assumptions, they are instructive for clarifying the conditions under which natural effects can 
be identified, and we take inspiration from this strand of literature in our approach in 
later sections. Finally, there have been some attempts to relax the strong independence tied 
to natural effects~\cite{robins2003,petersenetal2006}. 

The other strand of the literature has focused on developing theory for the identification and 
estimation of interventional effects. We cover the literature on the extension of interventional 
effects with multiple mediators in Section~\ref{chapter:interventional_effects} in 
greater detail. What we do not focus on is the recent literature on efficient semiparametric 
estimation of interventional 
effects for one and multiple mediator 
problems~\cite{rudolphetal2018,diazetal2021,benkeserran2021,rudolphetal2024}, which has been one 
of the most active areas of the mediation analysis literature in recent years. In general, a number
of other parametric strategies for estimating interventional effects have been proposed, including 
weighting~\citep{vanderweeleetal2014}, regression~\citep{vansteelandtdaniel2017}, and simulation 
methods~\citep{vansteelandtdaniel2017}. 

The interventional effects literature also has natural connections to the literature on stochastic
interventions~\citep{vanderweelerobins2012} and potential 
outcomes~\citep{gelmanmikhaeil2025}. Also in the tradition of population-level 
interventions,~\cite{hubbardvanderlaan2008} developed a framework for estimating the causal 
effects of potentially stochastic population-level interventions. A problem where such 
effect definitions might be useful would be a problem where we are interested in the health effects 
of a blanket smoking ban, but the health effects of such a policy themselves depend on the 
pre-existing distribution of smoking in the population.~\cite{diazvanderlaan2012} extended this 
framework to problems where the intervention is stochastic (for example, a policy to deter rather 
than to ban smoking). There is an interesting potential direct connection between these effects
and the VanderWeele-Robinson interventional effects we discuss in Section~\ref{section:vr_int_eff}, 
and it would be interesting to explore these connections further in future work. 

